\chapter{男性或女性的形成}\label{ch:g11:becoming-male-female}

人类胚胎在受精后第七周时，拥有一对既非睾丸也非卵巢的性腺，
以及并排的两套完整导管——一套可发育为男性生殖道，
另一套可发育为女性生殖道。在外观上，无法区分
将来发育为男孩的胚胎与将来发育为女孩的胚胎。
此后几周内，一套导管生长，另一套萎缩；
性腺分化为睾丸或卵巢，
到第四个月时，解剖结构已经确定。
这一决定由一个\omterm{def:g10:universal-dna:gene}{基因}作出，再由两种激素付诸实施。
本章沿着这一过程展开，
从受精时获得的\omterm{def:g11:cell-cycle-mitosis:chromosome}{染色体}讲到\omterm{prop:g11:becoming-male-female:puberty}{青春期}的变化。

\section{性染色体}

\begin{definition}[性染色体]\label{def:g11:becoming-male-female:chromosomes}
人类的 23 对\omterm{def:g11:cell-cycle-mitosis:chromosome}{染色体}中，有 22 对在两性中相同；
第 23 对由\emph{性染色体}\index{性染色体}组成：
女性为两条 X \omterm{def:g11:cell-cycle-mitosis:chromosome}{染色体}（XX），男性为一条 X 和一条 Y \omterm{def:g11:cell-cycle-mitosis:chromosome}{染色体}（XY）。
Y \omterm{def:g11:cell-cycle-mitosis:chromosome}{染色体}较小，携带的\omterm{def:g10:universal-dna:gene}{基因}很少，且只从父亲传给儿子；
每个卵\omterm{def:g10:cells-common-unit:cell}{细胞}都携带一条 X，精子则携带 X 或 Y 之一，
因此由父方的精子决定孩子的\omterm{def:g11:cell-cycle-mitosis:chromosome}{染色体}性别，
两种情况概率均等。
\end{definition}

\begin{proposition}[未分化阶段]\label{prop:g11:becoming-male-female:undifferentiated}
在第七周之前，两种\omterm{def:g11:cell-cycle-mitosis:chromosome}{染色体}性别的胚胎
生殖器官结构完全相同：一对未分化的性腺，
每个性腺旁各有两条导管——\emph{沃尔夫管}，
即男性生殖道的前体（附睾、输精管），以及
\emph{米勒管}，即女性生殖道的前体（输卵管、
子宫）。性发育是将这一共同的
结构转变为两种形式之一，
而启动这一转变的正是性腺的分化。
\end{proposition}
\admitted

\begin{omfigure}
\begin{tikzpicture}[font=\small, >=Stealth, line cap=round]
  % undifferentiated stage
  \begin{scope}
    \draw[thick, fill=gray!25] (0,2.2) ellipse (0.35 and 0.55);
    \node[font=\scriptsize, anchor=south] at (0,2.85) {性腺};
    \draw[very thick, omDef] (-0.25,1.6) -- (-0.25,0);
    \draw[very thick, omThm] (0.25,1.6) -- (0.25,0);
    \node[font=\scriptsize, omDef, anchor=east] at (-0.35,0.8) {沃尔夫管};
    \node[font=\scriptsize, omThm, anchor=west] at (0.35,0.45) {米勒管};
    \node[anchor=north, align=center] at (0,-0.2) {第 7 周：\\ 两套导管};
  \end{scope}
  \draw[->, thick] (1.6,1.6) -- (3.0,3.4) node[midway, above left, font=\scriptsize] {XY};
  \draw[->, thick] (1.6,0.9) -- (3.0,-1.0) node[midway, below left, font=\scriptsize] {XX};
  % male
  \begin{scope}[xshift=4.2cm, yshift=2.0cm]
    \draw[thick, fill=omDef!45] (0,2.2) ellipse (0.35 and 0.55);
    \node[font=\scriptsize, anchor=south] at (0,2.85) {睾丸};
    \draw[very thick, omDef] (-0.25,1.6) -- (-0.25,0);
    \draw[very thick, omThm, dotted] (0.25,1.6) -- (0.25,1.0);
    \node[font=\scriptsize, anchor=west, align=left] at (0.6,1.9) {睾酮：沃尔夫管 \\ 保留并发育};
    \node[font=\scriptsize, anchor=west, align=left] at (0.6,0.9) {AMH：米勒管 \\ 退化};
    \node[anchor=north, align=center] at (0,-0.2) {男性生殖道};
  \end{scope}
  % female
  \begin{scope}[xshift=4.2cm, yshift=-2.8cm]
    \draw[thick, fill=omThm!35] (0,2.2) ellipse (0.35 and 0.55);
    \node[font=\scriptsize, anchor=south] at (0,2.85) {卵巢};
    \draw[very thick, omDef, dotted] (-0.25,1.6) -- (-0.25,1.0);
    \draw[very thick, omThm] (0.25,1.6) -- (0.25,0);
    \node[font=\scriptsize, anchor=west, align=left] at (0.6,1.9) {无睾酮：沃尔夫管 \\ 退化};
    \node[font=\scriptsize, anchor=west, align=left] at (0.6,0.9) {无 AMH：米勒管 \\ 发育为输卵管和子宫};
    \node[anchor=north, align=center] at (0,-0.2) {女性生殖道};
  \end{scope}
\end{tikzpicture}

{\small 从共同结构到两种生殖道。睾丸借助其分泌的两种激素，
保留沃尔夫管并消除米勒管；在缺少这两种激素时，
米勒管发育，而沃尔夫管
退化。}
\end{omfigure}

\section{决定性别的基因}

\begin{proposition}[SRY与性腺]\label{prop:g11:becoming-male-female:sry}
Y \omterm{def:g11:cell-cycle-mitosis:chromosome}{染色体}携带一个\omterm{def:g10:universal-dna:gene}{基因}\emph{SRY}\index{SRY 基因}，该\omterm{def:g10:universal-dna:gene}{基因}在
第七周未分化的性腺中表达。其\omterm{def:g11:gene-expression:protein}{蛋白质}
开启使性腺转化为\emph{睾丸}的\omterm{def:g10:universal-dna:gene}{基因}。在
缺少 \omterm{prop:g11:becoming-male-female:sry}{SRY} 时（即 XX 胚胎中），性腺会在数周后
发育为\emph{卵巢}。决定性腺类型的是
\omterm{prop:g11:becoming-male-female:sry}{SRY} 的存在，而非 X \omterm{def:g11:cell-cycle-mitosis:chromosome}{染色体}的数目。
\end{proposition}
\begin{proof}[证据]
少数 \omterm{prop:g11:becoming-male-female:sry}{SRY} \omterm{def:g10:universal-dna:gene}{基因}发生\omterm{def:g11:mutations:mutation}{突变}的 XY 个体
发育出卵巢和女性身体；少数 XX 个体携带一段
带有 \omterm{prop:g11:becoming-male-female:sry}{SRY} 的 Y \omterm{def:g11:cell-cycle-mitosis:chromosome}{染色体}片段，该片段转移到 X \omterm{def:g11:cell-cycle-mitosis:chromosome}{染色体}上，发育出睾丸和男性身体。
小鼠实验：将小鼠的 \emph{Sry} \omterm{def:g10:universal-dna:gene}{基因}导入 XX 胚胎后，
该胚胎发育为雄性。Y \omterm{def:g11:cell-cycle-mitosis:chromosome}{染色体}上的一个\omterm{def:g10:universal-dna:gene}{基因}
既足以、也是必需的，使性腺形成睾丸。
\end{proof}

\begin{proposition}[胎儿睾丸的激素]\label{prop:g11:becoming-male-female:hormones}
睾丸形成后分泌两种激素。\emph{睾酮}\index{睾酮}
维持沃尔夫管，使其发育为男性内生殖道，
并使外生殖器发育为男性形态。
\emph{抗米勒管激素}（AMH）使米勒管
退化。此阶段卵巢两者均不分泌：缺少睾酮时
沃尔夫管退化，缺少 AMH 时米勒管
发育为输卵管、子宫和阴道上段。女性生殖道
因此，女性生殖道就是导管在缺乏
睾丸激素时所呈现的形式。
\end{proposition}
\begin{proof}[证据]
若斯特（1947）在分化之前摘除了子宫内兔胎儿的性腺：
结果每个胎儿，无论 XX 还是 XY，都发育出女性生殖道。
将一个睾丸移植到去势胎儿的性腺旁，
产生了男性生殖道；仅植入一块睾酮晶体，
保留了沃尔夫管，却未消除米勒管，后者
与之并存。因此，睾丸产生两种产物：一种构建男性导管
（睾酮），另一种（后被鉴定为 AMH）消除女性导管——而卵巢对女性生殖道的形成
并非
必需。
\end{proof}

\begin{omfigure}
\begin{tikzpicture}[font=\small, >=Stealth,
  b/.style={draw, rounded corners, align=center, text width=3.1cm, minimum height=1.1cm}]
  \node[b, fill=omDef!12] (a1) at (0,0) {胎儿，性腺已摘除};
  \node[b, fill=omThm!12] (a2) at (4.2,0) {无睾酮，\\ 无 AMH};
  \node[b, fill=white] (a3) at (8.4,0) {雌性生殖道，\\ 不论染色体如何};
  \draw[->] (a1) -- (a2); \draw[->] (a2) -- (a3);
  \node[b, fill=omDef!12] (b1) at (0,-1.7) {性腺已摘除，\\ 植入睾丸};
  \node[b, fill=omThm!12] (b2) at (4.2,-1.7) {睾酮 \\ 与 AMH};
  \node[b, fill=white] (b3) at (8.4,-1.7) {雄性生殖道};
  \draw[->] (b1) -- (b2); \draw[->] (b2) -- (b3);
  \node[b, fill=omDef!12] (c1) at (0,-3.4) {性腺已摘除，\\ 植入睾酮};
  \node[b, fill=omThm!12] (c2) at (4.2,-3.4) {睾酮，\\ 无 AMH};
  \node[b, fill=white] (c3) at (8.4,-3.4) {保留雄性导管；\\ 雌性导管同样保留};
  \draw[->] (c1) -- (c2); \draw[->] (c2) -- (c3);
\end{tikzpicture}

{\small 若斯特在兔胎儿上所做的三个实验及其解读。
雌性生殖道无需激素；雄性生殖道需要睾酮来构建，
还需要第二种睾丸激素 AMH 来消除
雌性导管。}
\end{omfigure}

\begin{example}[依次发生的级联过程]\label{ex:g11:becoming-male-female:cascade}
携带 Y 的精子受精；第 7 周，\omterm{prop:g11:becoming-male-female:sry}{SRY} 表达，性腺
发育为睾丸；第 8 至 12 周，睾酮与 AMH 作用：沃尔夫管
保留，米勒管消失，外生殖器形成男性形态；随后
进入静止期，直至\omterm{prop:g11:becoming-male-female:puberty}{青春期}。携带 X 的精子受精：无
\omterm{prop:g11:becoming-male-female:sry}{SRY}；第 10 至 12 周，性腺发育为卵巢；无睾丸
激素：米勒管发育，沃尔夫管退化，外生殖器
呈女性形态。每一步都依赖于前一步，
任一步出错都会改变其后的全部过程。
\end{example}

\section{青春期}

\begin{proposition}[青春期]\label{prop:g11:becoming-male-female:puberty}
生殖器官在出生前形成，童年时期一直处于
静止状态。在\emph{青春期}\index{青春期}（约 11 至 13 岁，女孩更早）时，
大脑中的一个中心开始
脉冲式释放一种激素，促使垂体分泌
两种激素 LH 与 FSH 进入血液。它们作用于性腺：
睾丸开始产生精子，并以成年水平分泌睾酮；
卵巢开始其周期性活动，并分泌雌激素。
性激素随后产生\emph{第二性征
}：包括生殖器官发育、体毛、
声音、乳房、肌肉与脂肪的分布、身高突增，
随后是骨骺闭合。
个体具备生殖能力，即\omterm{prop:g11:becoming-male-female:puberty}{青春期}完成。
\end{proposition}
\admitted

\begin{omfigure}
\begin{tikzpicture}
\begin{axis}[
  width=0.86\linewidth, height=5.6cm,
  axis x line=bottom, axis y line=left,
  xmin=0, xmax=20, ymin=0, ymax=110,
  xlabel={年龄（岁）}, ylabel={血液中性激素（占成年水平的 \%）},
  xlabel style={font=\small, at={(axis description cs:0.5,-0.12)}, anchor=north},
  ylabel style={font=\small},
  tick label style={font=\small}, xtick={0,4,8,12,16,20}, ytick={0,50,100},
  legend style={at={(0.03,0.95)}, anchor=north west, font=\small, draw=none},
]
\addplot[thick, omDef, smooth] coordinates {(0,3) (4,3) (8,3) (10,5) (11,12) (12,30) (13,55) (14,80) (15,95) (17,100) (20,100)};
\addplot[thick, omThm, smooth] coordinates {(0,3) (4,3) (8,4) (9,8) (10,20) (11,45) (12,70) (13,90) (15,100) (20,100)};
\legend{睾酮（男孩）, 雌激素（女孩）}
\end{axis}
\end{tikzpicture}

{\small 童年与\omterm{prop:g11:becoming-male-female:puberty}{青春期}血液中的性激素水平，
以占成年水平的百分比表示（近似平均值；个体时间差异
可达两三年）。在大脑发出信号之前，二者
均接近零。}
\end{omfigure}

\begin{example}[读懂曲线]\label{ex:g11:becoming-male-female:curves}
女孩的雌激素约从 9 岁开始上升，男孩的睾酮约从 11 岁开始上升；
两者在开始上升后两年内达到成年水平的一半，
再过两三年达到平台期。身高突增紧随上升之后约一年出现，
而骨骼停止生长则发生在激素达到成年水平之时——
这也是女孩开始得更早，因此平均而言停止生长也更早、
最终平均身高更矮的原因。
\end{example}

\section{性别的几层含义}

\begin{proposition}[区分各个层次]\label{prop:g11:becoming-male-female:levels}
“男性”与“女性”可用于若干层次，它们通常一致，
但彼此不同：
\begin{itemize}
  \item \emph{\omterm{def:g11:cell-cycle-mitosis:chromosome}{染色体}}性别：XX 或 XY，受精时即已确定；
  \item \emph{性腺}性别：卵巢或睾丸，由 \omterm{prop:g11:becoming-male-female:sry}{SRY} 决定；
  \item \emph{解剖}性别：生殖道与外生殖器，
        由胎儿性腺分泌的激素塑造，以及
        第二性征（在\omterm{prop:g11:becoming-male-female:puberty}{青春期}塑造）；
  \item \emph{性别认同}，即一个人认为自己是男人、
        女人，或都不是的感受，以及\emph{性取向}，即一个人被
        哪一类人吸引——这两者都是关于个人的事实，
        无法由前述层次推出，也不是人所能选择的。
\end{itemize}
本章的生物学内容描述前三个层次，它既不决定，
也不评判后者。
\end{proposition}
\admitted

\begin{example}[层次不一致的情形]\label{ex:g11:becoming-male-female:disagree}
约每几千人中有一人天生具有一种组合，
这种组合通常不会由上述级联过程产生：一名 XY 个体，其\omterm{def:g10:cells-common-unit:cell}{细胞}缺乏
睾酮受体，发育为女性体态，睾丸未下降，
且无子宫；一名 XX 个体，携带能形成睾丸的
\omterm{prop:g11:becoming-male-female:sry}{SRY} 易位，发育为男性体态；一名 XY 个体，若 \omterm{prop:g11:becoming-male-female:sry}{SRY}
无功能，则发育为女性。这些变异比常规情形更清楚地表明，
每一层次都是某一具体步骤的结果——
一个\omterm{def:g10:universal-dna:gene}{基因}、一种激素或一个受体——而且各层次
是相互独立的事实。
\end{example}

\begin{method}[根据具体情形推理]\label{met:g11:becoming-male-female:case}
\begin{enumerate}
  \item \omterm{def:g11:cell-cycle-mitosis:chromosome}{染色体}：\omterm{prop:g11:becoming-male-female:sry}{SRY} 是否存在且有功能？
  \item 性腺：若 \omterm{prop:g11:becoming-male-female:sry}{SRY} 起作用则为睾丸，否则为卵巢。
  \item 胎儿性腺分泌的激素：睾酮和 AMH，或二者皆无；
        组织能否作出反应（受体）？
  \item 导管：若两种激素均起作用，则保留中肾管、退化米勒管；
        若二者均不起作用，则发育米勒管；若仅睾酮起作用，
        而无 AMH，则两种导管均发育。
  \item \omterm{prop:g11:becoming-male-female:puberty}{青春期}：性腺分泌哪种性激素，决定
        第二性征——同样通过受体实现。
\end{enumerate}
\end{method}

\begin{remark}[一个基因，以及其后的一切]\label{rem:g11:becoming-male-female:onegene}
两种生殖器官结构的全部差异，可追溯到
胚胎第七周时，某一条\omterm{def:g11:cell-cycle-mitosis:chromosome}{染色体}上的某一\omterm{def:g10:universal-dna:gene}{基因}是否存在。此后的
过程是一连串普通的\omterm{def:g10:cells-common-unit:cell}{细胞}生物学——\omterm{def:g10:universal-dna:gene}{基因}表达、
激素、受体——即本学年每一章所述的内容。这一级联过程的起点
对两种性别是相同的；
而当睾丸不介入时，
它所构建的便是雌性形态。
\end{remark}

\section{练习}

\begin{exercise}[$\star$]\label{exo:g11:becoming-male-female:1}
子代的\omterm{def:g11:cell-cycle-mitosis:chromosome}{染色体}性别由哪一方的配子决定，
以及两种可能性为何相等？
\end{exercise}

\begin{exercise}[$\star$]\label{exo:g11:becoming-male-female:2}
描述七周龄胚胎的生殖器官结构。
\end{exercise}

\begin{exercise}[$\star$]\label{exo:g11:becoming-male-female:3}
\omterm{prop:g11:becoming-male-female:sry}{SRY} \omterm{def:g10:universal-dna:gene}{基因}起什么作用？它位于何处？
\end{exercise}

\begin{exercise}[$\star$]\label{exo:g11:becoming-male-female:4}
说出胎儿睾丸分泌的两种激素，
并说明各自对导管的作用。
\end{exercise}

\begin{exercise}[$\star$]\label{exo:g11:becoming-male-female:5}
什么触发\omterm{prop:g11:becoming-male-female:puberty}{青春期}？
第二性征有哪些？
\end{exercise}

\begin{exercise}[$\star\star$]\label{exo:g11:becoming-male-female:6}
在约斯特的实验中，被去势的 XY 胎儿表现出什么？
这说明了什么？
\end{exercise}

\begin{exercise}[$\star\star$]\label{exo:g11:becoming-male-female:7}
为什么仅植入睾酮时两套导管都保留了？
这说明了什么？
\end{exercise}

\begin{exercise}[$\star\star$]\label{exo:g11:becoming-male-female:8}
通过\omterm{prop:g10:universal-dna:universal}{转基因技术}，使一个 XX 小鼠胚胎获得 \emph{Sry} \omterm{def:g10:universal-dna:gene}{基因}。
预测其性腺、导管和外部生殖器形态，
并说明它成年后不能做什么。
\end{exercise}

\begin{exercise}[$\star\star$]\label{exo:g11:becoming-male-female:9}
根据激素图，各激素在什么年龄达到成年值的一半？
各自何时进入平台期？
\end{exercise}

\begin{exercise}[$\star\star$]\label{exo:g11:becoming-male-female:10}
利用曲线解释女孩为何比男孩更早停止生长。
\end{exercise}

\begin{exercise}[$\star\star$]\label{exo:g11:becoming-male-female:11}
一个 XY 胚胎携带某一\omterm{def:g11:mutations:mutation}{突变}，该\omterm{def:g11:mutations:mutation}{突变}使 \omterm{prop:g11:becoming-male-female:sry}{SRY} 蛋白失去功能。
沿级联过程推断，描述此人出生时及\omterm{prop:g11:becoming-male-female:puberty}{青春期}的情况。
\end{exercise}

\begin{exercise}[$\star\star\star$]\label{exo:g11:becoming-male-female:12}
一个 XY 胎儿的睾丸功能正常，但其\omterm{def:g10:cells-common-unit:cell}{细胞}缺乏
睾酮受体；AMH 作用正常。预测其性腺、两套导管、
外部生殖器形态，以及\omterm{prop:g11:becoming-male-female:puberty}{青春期}时的第二性征
（雌激素部分由睾酮转化而来）。
\end{exercise}

\begin{exercise}[$\star\star\star$]\label{exo:g11:becoming-male-female:13}
在某些爬行动物中，决定性别的是卵的温度，而非\omterm{def:g11:cell-cycle-mitosis:chromosome}{染色体}。
试说明性腺之后的级联过程在何种意义上可以与哺乳动物相同，
以及哪一步有所不同。
\end{exercise}

\begin{exercise}[$\star\star\star$]\label{exo:g11:becoming-male-female:14}
Y \omterm{def:g11:cell-cycle-mitosis:chromosome}{染色体}携带的\omterm{def:g10:universal-dna:gene}{基因}不足一百个，而 X \omterm{def:g11:cell-cycle-mitosis:chromosome}{染色体}约有一千个。
试解释为何一个人至少需要一条 X \omterm{def:g11:cell-cycle-mitosis:chromosome}{染色体}，却可以没有 Y \omterm{def:g11:cell-cycle-mitosis:chromosome}{染色体}，
以及为何 Y 连锁性状只能由父亲传给儿子。
\end{exercise}

\begin{exercise}[$\star\star\star$]\label{exo:g11:becoming-male-female:15}
试解释“雌性形态是发育的默认状态”这一说法
为何是关于激素的陈述，而非关于重要性的陈述，
并说明它对两种性别由共同方案演化而来
这一点有何启示。
\end{exercise}

\section{问题：约斯特的兔子}

\begin{problem}[{周末问题——发现睾丸两种激素的实验，
逐步重构其过程，然后应用于
三个人类案例以及青春期的时间}]
\label{pb:g11:becoming-male-female:1}
在兔子中，性腺约在妊娠第 19 天分化，导管约在
第 22 天分化。约斯特于子宫内对胎儿实施手术，手术日为
第 19 天，并在出生时（第 30 天）检查它们。

\medskip
\noindent\textbf{第一部分——实验。}
\begin{enumerate}
  \item 两种\omterm{def:g11:cell-cycle-mitosis:chromosome}{染色体}性别的胎儿，于第 19 天切除性腺：
        出生时均带有输卵管和子宫，且无雄性
        导管。这说明卵巢的作用，以及
        睾丸的作用，分别是什么？
  \item 为何必须在第 22 天之前手术？
  \item 第 19 天去势的胎儿，并将一个睾丸移植入
        腹腔：具有雄性导管，而无雌性导管。移植物提供什么，
        以及它如何到达导管？
  \item 第 19 天去势的胎儿，植入睾酮晶体，
        具有雄性导管，雌性导管也同时存在。
        睾酮起什么作用，又不起什么作用？
  \item 由此推断睾丸必定产生第二种物质，并说明
        其作用。之后如何证实它的存在？
\end{enumerate}

\medskip
\noindent\textbf{第二部分——切除一个性腺。}
\begin{enumerate}[resume]
  \item 于第 19 天从 XY 胎儿中切除一个睾丸。预测
        身体两侧的导管情况，并已知激素
        主要作用于较近的一侧。
  \item 将一个睾丸移植到一个 XX 胎儿的一侧卵巢旁。预测
        两侧的情况。
  \item 将一块睾酮晶体植入一个
        XX 胎儿的一侧卵巢旁。预测两侧的情况。
  \item 这些单侧实验为何能加强
        第一部分的结论？
  \item 解释为何移植了睾丸的 XX 胎儿
        成年后仍然不育。
\end{enumerate}

\medskip
\noindent\textbf{第三部分——三个人类案例。}
对每个案例应用\cref{met:g11:becoming-male-female:case}。
\begin{enumerate}[resume]
  \item 一个 XX 个体的一条 X \omterm{def:g11:cell-cycle-mitosis:chromosome}{染色体}上带有
        一段含 \omterm{prop:g11:becoming-male-female:sry}{SRY} 的 Y \omterm{def:g11:cell-cycle-mitosis:chromosome}{染色体}片段。
  \item 一个 XY 个体携带一个无功能的 \omterm{prop:g11:becoming-male-female:sry}{SRY}。
  \item 一个 XY 个体具有正常睾丸，但其\omterm{def:g10:cells-common-unit:cell}{细胞}无法
        对 AMH 产生反应。
  \item 三种情况中，哪一种个体可能具有生育能力？
        并说明理由。
  \item 三个案例说明\omterm{def:g11:cell-cycle-mitosis:chromosome}{染色体}性别与解剖性别之间
        存在怎样的关系？
\end{enumerate}

\medskip
\noindent\textbf{第四部分——数字中的\omterm{prop:g11:becoming-male-female:puberty}{青春期}。}
依据激素图，取生长突增期于激素达到成人水平的
20\% 后一年开始，骨骼则在激素达到
100\% 时停止生长。
\begin{enumerate}[resume]
  \item 估算女孩和男孩生长突增期开始的年龄，
        以及生长停止的年龄。
  \item 两性各自从突增期开始到结束，
        经历了多少年的生长？
  \item 男孩在突增期每年长高 \qty{8}{cm}，女孩则为
        \qty{7}{cm}。估算两者在突增期间各长高多少，
        并评论两性成年身高的平均差异
        （约 \qty{13}{cm}）。
  \item 若女孩的雌激素升高始于 7 岁，她将早期快速生长，
        并早早停止生长。解释医生为何有时用药物
        延缓过早的\omterm{prop:g11:becoming-male-female:puberty}{青春期}。
  \item 归纳结论：约斯特的兔子实验揭示了两种激素，
        各自的作用，以及开启产生它们的器官的
        那一个\omterm{def:g10:universal-dna:gene}{基因}。
\end{enumerate}
\end{problem}
